\documentclass[12pt]{article}
\usepackage{amsmath, amssymb, amsthm}
\usepackage{geometry}
\geometry{margin=1in}
\usepackage{mathrsfs}
\usepackage{bm}
\usepackage[utf8]{inputenc}
\usepackage{hyperref}
\usepackage{listings}
\newtheorem{theorem}{Theorem}
\newtheorem{lemma}[theorem]{Lemma}
\newtheorem{corollary}[theorem]{Corollary}
\newtheorem{definition}[theorem]{Definition}
\newtheorem{conjecture}[theorem]{Conjecture}
\newtheorem{proposition}[theorem]{Proposition}

\newcommand{\RR}{\mathbb{R}}
\newcommand{\ZZ}{\mathbb{Z}}
\newcommand{\QQ}{\mathbb{Q}}
\newcommand{\EE}{\mathbb{E}}
\newcommand{\A}{\mathcal{A}}
\newcommand{\G}{\mathcal{G}}
\newcommand{\ResGraph}{\mathbf{ResGraph}}
\newcommand{\ResGraphZero}{\mathbf{ResGraph}_0}
\newcommand{\Weyl}{W}
\newcommand{\Eeight}{E_8}
\newcommand{\PhiAdj}{\Phi}
\newcommand{\clique}{\text{clique}}

\begin{document}

\begin{titlepage}
    \centering
 
    
  \LARGE\textbf{Atlas Embeddings: Exceptional Lie Groups from First Principles}
   \vspace*{0.5em}
  {\Large\bfseries \\Mathematical Foundations and Main Theorems \par}
   \vspace{1cm}
     \Large{UOR FOUNDATION} \\
  \vspace{1cm}
  \text{Ryan O. Van Gelder} \\ \vspace{1cm} 
    {\normalsize October 2025 \par}
    \vfill

   
    \vspace{0.5em}
    \begin{minipage}{0.85\textwidth}
        \small
       We present a publication-ready specification of the $E_8$ root system with a deterministic paste-stable ordering, an injective reconstructive labeling, and a category of \emph{ResGraph} objects capturing induced substructures and symmetries. Simple roots are given in a Bourbaki-consistent realization (trivalent at node $4$) with inner product $\langle\alpha_i,\alpha_j\rangle=-1$ on Dynkin edges and $2$ on the diagonal.

A $\Sigma$-calculus encodes constructions via three primitives: parabolic restriction $\mathrm{Par}_I$, admissible fold $\mathrm{fold}_\Gamma$ (including quotient folds), and junction tagging $\mathrm{Aug}_J$. Using these, we supply five paste-stable $\Sigma$-terms whose reflection-closures recover the classical types inside $E_8$: $G_2$ from a $D_4$ triality fold, $F_4$ from an $E_6$ quotient fold, $E_6$ from a standard parabolic, $E_7$ from a parabolic with admissible augmentation, and $E_8$ from the Atlas simple system itself.

Each construction is backed by machine checks: Cartan recovery, determinant and rank, integrality of Cartan integers on the closure, Coxeter parameters, root counts, and Weyl orders. Verified outputs are: $G_2$ (12 roots, $|W|=12$), $F_4$ (48 roots, $|W|=1152$), $E_6$ (72 roots, $|W|=51{,}840$), $E_7$ (126 roots, $|W|=2{,}903{,}040$), and $E_8$ (240 roots, $|W|=696{,}729{,}600$). The framework fixes conventions end-to-end, yields reproducible code paths, and isolates admissibility conditions for folds and augmentations. This provides a minimal, verifiable, and implementation-ready atlas for $E_8$ and its induced subsystems.


    \end{minipage}

\end{titlepage}
\clearpage
\clearpage

\newpage
\begin{multicols}{2}
\begin{singlespace}
\tableofcontents
\end{singlespace}
\end{multicols}

\thispagestyle{empty}
\newpage

\textbf{Atlas of Resonance Classes}.


\section{The Atlas of Resonance Classes}

\paragraph{Object.} A canonical 96–element set of resonance classes \(\RninetySix=\{v_1,\dots,v_{96}\}\subset \Lambda_{\E}\) with a fixed adjacency predicate
\[
(i\sim j) \iff \ip{v_i}{v_j} = \AdjTarget, \quad \AdjTarget\in\{-1,+1\}.
\]
The two sign choices are Weyl–equivariant; pick one and use it consistently.


\paragraph{Multiplicity profile.} \RninetySix is split into two resonance levels with counts
\[
32 \;\text{elements at level}\; 2,\qquad 64 \;\text{elements at level}\; 3.
\]
These labels are intrinsic to the atlas and do not rely on any particular coordinate chart.


\paragraph{Embedding.} There exists an embedding \(\iota: \RninetySix\hookrightarrow \Lambda_{\E}\) unique up to \(\W\). Any two embeddings differ by an element of \(\W\).


\paragraph{Graph structure.} Let \(G(\RninetySix)\) be the simple graph on 96 vertices with the above adjacency. Key invariants:
\begin{itemize}[nosep]
\item Degree partition: \(\{d_a^{\times 32},\; d_b^{\times 64}\}\) for some integers \(d_a\neq d_b\).
\item Induced exceptional subgraphs: counts \(N_{E_6}, N_{E_7}, N_{E_8}\) of \(E_6,E_7,E_8\) Dynkin diagrams.
\item Automorphisms: \(\operatorname{Aut}(G)\) contains the image of \(\W\) stabilizing \(\RninetySix\).
\end{itemize}


\paragraph{Verification sketch.} Given coordinates for \(\RninetySix\) in \(\Lambda_{\E}\):
\begin{enumerate}[nosep]
\item Use exact rationals to compute the matrix \(M_{ij}=\ip{v_i}{v_j}\).
\item Form the adjacency matrix \(A\) by testing \(M_{ij}=\AdjTarget\) for \(i\neq j\).
\item Check the 96–way degree partition equals \(\{d_a^{\times 32}, d_b^{\times 64}\}\).
\item Count induced \(E_6, E_7, E_8\) subgraphs as a sanity check of exceptional closures.
\item Record a certificate: the raw \(M\), the edge set, and hashes of both.
\end{enumerate}

\begin{definition}[Atlas Graph]
The \textbf{Atlas} $\A = (V_\A, E_\A, \lambda_\A, U_\A)$ is a 96-vertex graph with:
\begin{itemize}
\item $V_\A$: 96 vertices representing resonance classes
\item $E_\A \subseteq V_\A \times V_\A$: edges defined by the adjacency predicate
\item $\lambda_\A: V_\A \to R_{E_8}$: vertex labeling into $E_8$ roots
\item $U_\A \subseteq V_\A$: distinguished unity vertices
\end{itemize}
\end{definition}

\begin{definition}[Adjacency Predicate]
For $t \in \{0, \pm1\}$, define:
\[
\PhiAdj(t) \iff (t = -1)
\]
Two vertices $v, w \in V_\A$ are adjacent iff $\langle \lambda_\A(v), \lambda_\A(w) \rangle = -1$.
\end{definition}

\subsection{The Category of Resonance Graphs}

\begin{definition}[Resonance Graph Object]
A \textbf{ResGraph object} is a tuple $\G = (V, E, \lambda, U)$ where:
\begin{enumerate}
\item $V$ is a finite set of vertices
\item $E \subseteq \{\{v,w\} : v,w \in V, v \neq w\}$ is a set of edges
\item $\lambda: V \to R_{E_8}$ labels vertices with $E_8$ roots
\item $U \subseteq V$ is a distinguished unity subset
\item Adjacency rule: $\{v,w\} \in E$ iff $\PhiAdj(\langle \lambda(v), \lambda(w) \rangle)$
\end{enumerate}
\end{definition}

\begin{definition}[ResGraph Morphism]
A morphism $\varphi: (V,E,\lambda,U) \to (V',E',\lambda',U')$ is a function $\varphi: V \to V'$ such that:
\begin{enumerate}
\item (Graph homomorphism) $\{v,w\} \in E \Rightarrow \{\varphi(v), \varphi(w)\} \in E'$
\item (Label compatibility) $\exists g \in \Weyl(E_8)$ such that $\lambda'(\varphi(v)) = g \cdot \lambda(v)$ for all $v \in V$
\item (Unity preservation) $\varphi(U) \subseteq U'$
\end{enumerate}
\end{definition}

\subsection{Main Theorem: Atlas to $E_8$ Embedding}

\begin{theorem}[Atlas Embedding]
There exists an injective map $f: V_\A \to R_{E_8}$ such that:
\begin{enumerate}
\item For all $v,w \in V_\A$, $\{v,w\} \in E_\A$ iff $\langle f(v), f(w) \rangle = -1$
\item $|\operatorname{im}(f)| = 96$
\item (Uniqueness up to Weyl group) If $f'$ also satisfies (1) and (2), then there exists $w \in \Weyl(E_8)$ with $f' = w \circ f$
\end{enumerate}
\end{theorem}

\begin{proof}[Proof Outline]
\textbf{Existence:} Construct $f$ explicitly by providing coordinates for all 96 vertices in the canonical $E_8$ basis. Verify adjacency by checking all $\binom{96}{2}$ pairs satisfy $\PhiAdj$.

\textbf{Uniqueness:} 
\begin{enumerate}
\item Fix a seed vertex $v_0$ and normalize $f(v_0)$ to a chosen root via Weyl action
\item Neighbors of $v_0$ must map to the 56 roots with inner product $-1$ with $f(v_0)$
\item Local adjacency patterns force unique assignments via connectedness of $\A$
\item Any alternative embedding differs by a global Weyl transformation
\end{enumerate}
\end{proof}

\subsection{Categorical Operations in $\ResGraph$}

\begin{definition}[Reflection Closure]
For $\G = (V,E,\lambda,U)$, define $\operatorname{RC}(\G)$ as the smallest ResGraph containing $\G$ whose label set is closed under reflections:
\[
\Lambda_{k+1} = \Lambda_k \cup \{s_\alpha(\beta) : \alpha, \beta \in \Lambda_k\} \cap R_{E_8}
\]
with $\Lambda_0 = \lambda(V)$.
\end{definition}

\begin{definition}[Resonance Product]
For $\G_1, \G_2$, define the partial sum graph:
\[
\G_1 \boxtimes \G_2 = \left\{ (v_1,v_2) : \lambda_1(v_1) + \lambda_2(v_2) \in R_{E_8} \right\}/\sim
\]
where $(v_1,v_2) \sim (v_1',v_2')$ if their label sums coincide.
\end{definition}

\begin{definition}[Quotient by Symmetry]
For $H \leq \operatorname{Aut}(\G)$, define $\G/H$ with vertices as $H$-orbits, labels by representatives, and edges between orbits when any representatives are adjacent.
\end{definition}

\subsection{Exceptional Group Constructions}

\begin{theorem}[$G_2$ Construction]
The exceptional group $G_2$ emerges as:
\[
G_2 = \operatorname{RC}(\A \boxtimes C_3)/H
\]
where $C_3$ is a 3-element resonance graph and $H$ an appropriate symmetry group. The construction yields:
\begin{itemize}
\item 12 roots (6 short, 6 long)
\item Rank 2, Cartan matrix $\begin{bmatrix} 2 & -3 \\ -1 & 2 \end{bmatrix}$
\item Weyl group order 12
\end{itemize}
\end{theorem}

\begin{theorem}[$F_4$ Construction]
\[
F_4 = \operatorname{RC}(\A/\tau)
\]
where $\tau$ is the mirror symmetry involution, yielding:
\begin{itemize}
\item 48 roots (24 short, 24 long)
\item Rank 4
\item Weyl group order 1152
\end{itemize}
\end{theorem}

\begin{theorem}[$E_6$, $E_7$, $E_8$ Constructions]
\begin{itemize}
\item $E_6 = \operatorname{RC}(\A[\ell = c])$ via hyperplane filtration
\item $E_7 = \operatorname{Aug}(\A)$ via root addition closure  
\item $E_8 = \operatorname{RC}(\A)$ via full reflection closure
\end{itemize}
\end{theorem}

\subsection{Initiality and Universal Property}

\begin{conjecture}[Atlas Initiality]
The Atlas $\A$ is initial in $\ResGraph$: for every object $\G \in \ResGraph$, there exists a unique morphism $\varphi: \A \to \G$.
\end{conjecture}

\begin{definition}[Precursor Category]
Define $\ResGraphZero$ with objects $(V,E,U)$ (forgetting labels) and morphisms preserving graph structure and unity. There is a faithful functor $J: \ResGraph \to \ResGraphZero$.
\end{definition}

\begin{lemma}[Lifting Lemma]
Any $\ResGraphZero$-morphism $f: J(\A) \to J(\G)$ lifts uniquely to a $\ResGraph$-morphism $\tilde{f}: \A \to \G$ if $\G$ satisfies the $E_8$-liftability condition.
\end{lemma}

\subsection{Action Functional and Variational Foundation}

\begin{definition}[Flag-Resonance Complex]
For $\G = (V,E,\lambda,U)$, define the cellular complex $\mathsf{FR}(\G)$ by:
\begin{itemize}
\item 0-cells: $V$
\item 1-cells: $E$  
\item $k$-cells: $\clique_{k+1}(\G)$ (cliques of size $k+1$)
\end{itemize}
\end{definition}

\begin{definition}[Action Functional]
For $\varphi \in C^{k-1}(\mathsf{FR}(\G); A)$, define:
\[
S_k[\varphi] = \sum_{\sigma \in \mathsf{FR}^{(k)}} \varphi(\partial \sigma)
\]
The Atlas emerges as stationary configuration of $S_k$ for appropriate $k$ and coefficient group $A$.
\end{definition}

\subsection{Computational Verification Framework}

\begin{theorem}[Exact Arithmetic Verification]
All constructions and verifications use exact rational arithmetic:
\begin{enumerate}
\item Vertex coordinates in $\QQ$ or $\frac12\ZZ$
\item Inner products computed exactly via $\langle x,y \rangle = \sum x_i y_i$
\item Adjacency verified by $\PhiAdj(\langle \lambda(v), \lambda(w) \rangle)$
\item Cartan matrices derived from simple systems
\end{enumerate}
\end{theorem}

\begin{proposition}[Reproducible Artifacts]
The theory is computationally certified by:
\begin{itemize}
\item Complete vertex/edge listings for $\A$
\item Explicit coordinate tables for all exceptional groups
\item Verification scripts confirming all mathematical claims
\item Flag complex cell counts justifying the 12,288-cell structure
\end{itemize}
\end{proposition}


\section{The $E_8$ Root System}

The classification of simple Lie algebras reveals five exceptional cases: $G_2$, $F_4$, $E_6$, $E_7$, and $E_8$. While these have been extensively studied since their discovery by Killing and Cartan, their emergence from first principles remains mysterious. This work presents a unified construction from a single combinatorial object: the 

\begin{definition}[$E_8$ Root System]
The $E_8$ root system $R_{E_8} \subset \RR^8$ consists of 240 vectors of squared norm 2:
\begin{itemize}
\item 112 integer roots: $\pm e_i \pm e_j$ for $1 \leq i < j \leq 8$
\item 128 half-integer roots: $(\pm \tfrac12, \pm \tfrac12, \ldots, \pm \tfrac12)$ with an even number of minus signs
\end{itemize}
The inner product between distinct roots takes values in $\{0, \pm1\}$.
\end{definition}

\begin{definition}[Weyl Group]
The Weyl group $\Weyl(E_8)$ is generated by reflections $s_\alpha(x) = x - \frac{2\langle x, \alpha\rangle}{\langle \alpha, \alpha\rangle} \alpha$ for $\alpha \in R_{E_8}$.
\end{definition}

\subsection{Root System G$_2$ (product)}\label{sec:g2-fold}

\paragraph{$\Sigma$-term.}
Work in Bourbaki $E_8$ with simple roots $\{\alpha_i\}_{i=1}^8$ and $\langle\alpha_i,\alpha_j\rangle=-1$ exactly along Dynkin edges. Inside $E_8$, the subset $\{\alpha_1,\alpha_2,\alpha_3,\alpha_4\}$ is a $D_4$ with hub $\alpha_3$ and legs $\alpha_1,\alpha_2,\alpha_4$.
Define the \text{$\Sigma$-term}
\[
t_{G_2}\;:=\;\mathrm{fold}_{C_3}\!\big(\mathrm{Par}_{\{1,2,4\}};\,\mathrm{Aug}_{3}\big)(\text{Atlas marking}),
\]
so that
\[
\beta_\ell:=\alpha_1+\alpha_2+\alpha_4,\qquad \beta_s:=\alpha_3.
\]

\paragraph{Cartan check.}
With $\langle\alpha_i,\alpha_i\rangle=2$, legs pairwise orthogonal, and each leg meeting the hub once,
\[
\|\beta_\ell\|^2=6,\quad \|\beta_s\|^2=2,\quad \langle \beta_\ell,\beta_s\rangle=-3,
\]
hence
\[
A=\begin{pmatrix}
2 & \dfrac{2\langle\beta_\ell,\beta_s\rangle}{\|\beta_s\|^2}\\[4pt]
\dfrac{2\langle\beta_s,\beta_\ell\rangle}{\|\beta_\ell\|^2} & 2
\end{pmatrix}
=
\begin{pmatrix}2&-3\\-1&2\end{pmatrix},
\]
the crystallographic Cartan matrix of type $G_2$.

\paragraph{Angle and Coxeter parameter.}
\[
\cos\theta=\frac{\langle\beta_s,\beta_\ell\rangle}{\|\beta_s\|\,\|\beta_\ell\|}=-\frac{3}{\sqrt{12}}=-\frac{\sqrt{3}}{2}\;\Rightarrow\; \theta=150^\circ.
\]
For simple roots, $\theta=\pi-\frac{\pi}{m}$, so $m=\frac{\pi}{\pi-\theta}=6$.

\paragraph{Root strings.}
With $\alpha=\beta_s$ and $\beta=\beta_\ell$:
\[
\langle\beta,\alpha^\vee\rangle=-3\ \Rightarrow\ \text{$\alpha$-string through $\beta$ has length }1-(-3)=4,
\]
\[
\langle\alpha,\beta^\vee\rangle=-1\ \Rightarrow\ \text{$\beta$-string through $\alpha$ has length }1-(-1)=2,
\]
matching $G_2$ string lengths and ensuring crystallographic closure.

\begin{proposition}
The reflection-closure of $\{\beta_\ell,\beta_s\}$ in the $E_8$ ambient is the $G_2$ root system. The Weyl group is dihedral of order $12$.
\end{proposition}

\begin{proof}[Proof outline]
Triality $C_3$ permutes the three $D_4$ legs, and folding along $C_3$ is admissible in simply-laced type because nodes in the same orbit are nonadjacent. The Gram data above yield the $G_2$ Cartan matrix. The angle gives $m=6$, so the rank-$2$ Coxeter group generated by $s_{\beta_s},s_{\beta_\ell}$ is $D_{12}$ of order $12$. Root-string lengths $(4,2)$ confirm crystallographic closure with long:short norm ratio $3:1$. \qedhere
\end{proof}

\subsection{Root System F$_4$ (quotient fold of an $E_6\subset E_8$ slice)}\label{sec:f4-fold}

\paragraph{$\Sigma$-term.}
Let $\{\alpha_i\}_{i=1}^8$ be the Bourbaki $E_8$ simple roots with $\langle\alpha_i,\alpha_j\rangle=-1$ on Dynkin edges. Inside $E_8$, take a symmetric $E_6$ slice written as a chain with a central node and an attached arm:
\[
L-\ell-\boxed{c}-r-R,\qquad u\text{ attached to }c .
\]
Let $\Gamma=C_2$ act by $L\leftrightarrow R$ and $\ell\leftrightarrow r$, fixing $c,u$. Define the quotient (orbit–average) fold
\[
t_{F_4}\;:=\;\mathrm{fold}^{\mathrm{quot}}_{C_2}\!\big(\mathrm{Par}_{E_6};\,\mathrm{Aug}_{\{\text{center, arm}\}}\big)(\text{Atlas}),
\qquad
p(O)=\frac{1}{|O|}\sum_{\alpha\in O}\alpha .
\]

\paragraph{Folded generators.}
\[
\beta_1=\tfrac12(L+R),\quad
\beta_2=\tfrac12(\ell+r),\quad
\beta_3=c,\quad
\beta_4=u .
\]

\paragraph{Gram data in $E_8$.}
Lengths:
\(
\|\beta_1\|^2=\|\beta_2\|^2=1\) (short), \(
\|\beta_3\|^2=\|\beta_4\|^2=2\) (long).
Pairings:
\[
\langle\beta_1,\beta_2\rangle=-\tfrac12,\qquad
\langle\beta_2,\beta_3\rangle=-1,\qquad
\langle\beta_3,\beta_4\rangle=-1,
\]
others $0$.

\paragraph{Cartan matrix.}
With $A_{ij}=2\langle\beta_i,\beta_j\rangle/\|\beta_j\|^2$,
\[
A=
\begin{pmatrix}
2&-1&0&0\\
-1&2&-1&0\\
0&-2&2&-1\\
0&0&-1&2
\end{pmatrix},
\]
which is the $F_4$ Cartan (double bond between nodes $2$ and $3$, arrow toward the short root).

\paragraph{Angles and strings.}
$m_{12}=3$ (short–short, $\cos120^\circ$), $m_{23}=4$ (short–long), $m_{34}=3$ (long–long). Root strings across $(\beta_2,\beta_3)$ have lengths $2$ and $3$ since $\langle\beta_2,\beta_3^\vee\rangle=-1$ and $\langle\beta_3,\beta_2^\vee\rangle=-2$.

\begin{proposition}
The reflection-closure of $\{\beta_1,\beta_2,\beta_3,\beta_4\}$ is the $F_4$ root system. The Weyl group has order $1152$.
\end{proposition}

\begin{proof}[Proof outline]
The $C_2$ involution swaps two symmetric arms of the $E_6$ slice, so the quotient fold via orbit-averaging is admissible. The Gram data yield the $F_4$ Cartan above. The adjacent Coxeter parameters $(3,4,3)$ match $F_4$. Root-string lengths $(2,3)$ across the double bond certify crystallographic closure with the expected short/long structure. Therefore the reflection-closure is $F_4$ and the Weyl group order is $1152$.
\end{proof}

\subsection{Root System E$_6$ (filtration)}\label{sec:e6-filtration}

\paragraph{$\Sigma$-term.}
Work in Bourbaki $E_8$ with $\langle\alpha_i,\alpha_j\rangle=-1$ on Dynkin edges. Select the standard $E_6$ parabolic and tag the central node for checks:
\[
t_{E_6}:=\mathrm{Par}_{I}\circ \mathrm{Aug}_{\{4\}}\;(\text{Atlas}),\qquad
I=\{1,2,3,4,5,6\}.
\]
As a filtration one may use
\[
E_8 \xrightarrow{\ \mathrm{Par}_{\{1,\dots,7\}}\ } E_7
\;\xrightarrow{\ \mathrm{Par}_{\{1,2,3,4,5,6\}}\ } E_6 .
\]

\paragraph{Simple system produced.}
Set $\beta_i:=\alpha_i$ for $i\in I$. Adjacencies (trivalent node at $4$):
\[
1\!-\!3,\quad 3\!-\!4,\quad 4\!-\!5,\quad 5\!-\!6,\quad 2\!-\!4,
\]
others $0$. All $\|\beta_i\|^2=2$.

\paragraph{Cartan matrix.}
With $A_{ij}=2\langle\beta_i,\beta_j\rangle/\|\beta_j\|^2$,
\[
A=
\begin{pmatrix}
2&0&-1&0&0&0\\
0&2&0&-1&0&0\\
-1&0&2&-1&0&0\\
0&-1&-1&2&-1&0\\
0&0&0&-1&2&-1\\
0&0&0&0&-1&2
\end{pmatrix},
\qquad \det A=3,\ \mathrm{rank}=6.
\]

\paragraph{Reflection-closure and size.}
Since $A$ is simply-laced, crystallographic, and indecomposable, the reflection-closure is the $E_6$ root system with $|\Phi|=72$ and Coxeter number $h=12$.

\paragraph{Weyl order.}
Exponents $(1,4,5,7,8,11)$ give
\[
|W(E_6)|=\prod_{i=1}^6 (m_i+1)=2\cdot5\cdot6\cdot8\cdot9\cdot12=51{,}840.
\]

\subsection{Root System E$_7$ (augmentation)}\label{sec:e7-augmentation}

\paragraph{$\Sigma$-term.}
Work in Bourbaki $E_8$ with $\langle\alpha_i,\alpha_j\rangle=-1$ on Dynkin edges. Select the standard $E_7$ parabolic and tag the junction (node $4$) for checks:
\[
t_{E_7}:=\mathrm{Aug}_{\{4\}}\circ \mathrm{Par}_{I}\;(\text{Atlas}),\qquad
I=\{1,2,3,4,5,6,7\}.
\]
The augmentation is admissible: it adds no edges and preserves integrality.

\paragraph{Simple system produced.}
Set $\beta_i:=\alpha_i$ for $i\in I$. Adjacencies (trivalent at $4$):
\[
1\!-\!3,\quad 3\!-\!4,\quad 4\!-\!5,\quad 5\!-\!6,\quad 6\!-\!7,\quad 2\!-\!4,
\]
others $0$. All $\|\beta_i\|^2=2$.

\paragraph{Cartan matrix.}
With $A_{ij}=2\langle\beta_i,\beta_j\rangle/\|\beta_j\|^2$,
\[
A=
\begin{pmatrix}
2&0&-1&0&0&0&0\\
0&2&0&-1&0&0&0\\
-1&0&2&-1&0&0&0\\
0&-1&-1&2&-1&0&0\\
0&0&0&-1&2&-1&0\\
0&0&0&0&-1&2&-1\\
0&0&0&0&0&-1&2
\end{pmatrix},
\qquad \det A=2,\ \mathrm{rank}=7.
\]

\paragraph{Reflection-closure and size.}
Since $A$ is simply-laced, crystallographic, indecomposable, the reflection-closure is the $E_7$ root system with $|\Phi|=126$ and Coxeter number $h=18$.

\paragraph{Weyl order.}
Exponents $(1,5,7,9,11,13,17)$ give
\[
|W(E_7)|=\prod_{i=1}^7 (m_i+1)=2\cdot6\cdot8\cdot10\cdot12\cdot14\cdot18
=2{,}903{,}040.
\]


Work in $\mathbb{R}^8$ with the standard dot product $\langle x,y\rangle=\sum_{i=1}^8 x_i y_i$.
Define the even unimodular lattice
\[
E_8=\{x\in\mathbb{Z}^8\}\ \cup\ \bigl\{x\in(\mathbb{Z}+\tfrac12)^8:\ \textstyle\sum_{i=1}^8 x_i\in 2\mathbb{Z}\bigr\}.
\]
The root set is
\[
R=\{\alpha\in E_8:\ \langle \alpha,\alpha\rangle=2\},\qquad |R|=240.
\]
Two disjoint types:
\begin{align*}
\text{Type I: }&\text{all permutations of }(\pm1,\pm1,0,\dots,0)\ \ \text{(no sign-parity restriction)}\quad\Rightarrow\ 112\ \text{roots},\\
\text{Type II: }&(\pm\tfrac12,\dots,\pm\tfrac12)\ \text{with even parity (even \# of $-1$ entries)}\quad\Rightarrow\ 128\ \text{roots}.
\end{align*}
Possible inner products among roots are $\{-2,-1,0,1,2\}$.


\subsection{Cartan Matrix, Simple Roots, and Weyl Group}
Adopt the Bourbaki $E_8$ Cartan matrix with the trivalent node at $3$:
\[
A=\begin{pmatrix}
2&0&-1&0&0&0&0&0\\
0&2&-1&0&0&0&0&0\\
-1&-1&2&-1&0&0&0&0\\
0&0&-1&2&-1&0&0&0\\
0&0&0&-1&2&-1&0&0\\
0&0&0&0&-1&2&-1&0\\
0&0&0&0&0&-1&2&-1\\
0&0&0&0&0&0&-1&2
\end{pmatrix}.
\]
A compatible simple system $\{\alpha_i\}_{i=1}^8\subset\mathbb{R}^8$ is:
\[
\begin{aligned}
\alpha_1&=\tfrac12(1,-1,-1,-1,-1,-1,-1,1),&
\alpha_2&=(1,1,0,0,0,0,0,0),\\
\alpha_3&=(-1,0,1,0,0,0,0,0),&
\alpha_4&=(0,0,-1,1,0,0,0,0),\\
\alpha_5&=(0,0,0,-1,1,0,0,0),&
\alpha_6&=(0,0,0,0,-1,1,0,0),\\
\alpha_7&=(0,0,0,0,0,-1,1,0),&
\alpha_8&=(0,0,0,0,0,0,-1,1).
\end{aligned}
\]
Then $\langle\alpha_i,\alpha_i\rangle=2$ and $\langle\alpha_i,\alpha_j\rangle=-1$ exactly on the edges of $A$; all other off-diagonals are $0$.
Define reflections
\[
s_i(x)=x-\frac{2\langle x,\alpha_i\rangle}{\langle\alpha_i,\alpha_i\rangle}\,\alpha_i
\quad(\langle\alpha_i,\alpha_i\rangle=2),
\]
and the Weyl group $W(E_8)=\langle s_1,\dots,s_8\rangle$, which acts transitively on $R$.

\subsection{Paste-Stable Ordering \texorpdfstring{$\iota$}{iota}}
We specify a deterministic total order $\iota:R\to\{1,\dots,240\}$:
\begin{enumerate}[itemsep=2pt,topsep=4pt,leftmargin=16pt]

\item \textbf{Type I ordering.} Write $\alpha=\epsilon_i e_i+\epsilon_j e_j$ with $i<j$ and $\epsilon_i,\epsilon_j\in\{\pm1\}$. Order by the tuple
\[
(i,\ j,\ \epsilon_i,\ \epsilon_j)\qquad\text{under the bit order }-1<+1.
\]
\item \textbf{Type II ordering.} Use the sign vector $\sigma(\alpha)\in\{\pm1\}^8$ (even parity). Order lexicographically on $(\sigma_1,\dots,\sigma_8)$ with $-1<+1$.
\item Do not identify $\alpha$ with $-\alpha$ for ordering.
\end{enumerate}

\subsection{Adjacency Predicate and Root Graph}
Define the $W(E_8)$-invariant adjacency
\[
\Phi(\alpha,\beta)\iff \langle \alpha,\beta\rangle=-1.
\]
This yields a simple undirected $56$-regular graph on vertex set $R$.

\subsection{Injective, Reconstructive Labels}
Define $\lambda:R\to \{\mathrm{I},\mathrm{II}\}\times\mathcal{D}$ by
\[
\lambda(\alpha)=
\begin{cases}
(\mathrm{I},(i,j),(\epsilon_i,\epsilon_j))& \text{if }\alpha=\epsilon_i e_i+\epsilon_j e_j,\ i<j,\\[2pt]
(\mathrm{II},\sigma(\alpha))& \text{if }\alpha\in\{\pm\tfrac12\}^8\text{ with even parity},
\end{cases}
\]
where $\sigma(\alpha)$ is the $8$-bit sign vector. The reconstruction map
\[
\mathrm{rec}\circ\lambda=\mathrm{id}_V\qquad\text{for any }V\subseteq R
\]
is immediate from the data.

\subsection{ResGraph Objects, Morphisms, and Atlas}
\subsubsection*{Objects}
A \emph{ResGraph} is a tuple $G=(V,E,\lambda,U)$ with
\begin{align*}
&V\subseteq R\ \text{specified by indices }I\subseteq\{1,\dots,240\}: &&V=\{\alpha\in R:\ \iota(\alpha)\in I\},\\
&E=\bigl\{\{\alpha,\beta\}\subseteq V:\ \Phi(\alpha,\beta)\bigr\},\qquad
&&\lambda=\text{labels above, restricted to }V,\\
&U=\{w|_V:\ w\in W(E_8),\ w(V)=V\}\subseteq \mathrm{Aut}(G).
\end{align*}

\subsection*{Morphisms}
A morphism $f:(V,E,\lambda,U)\to(V',E',\lambda',U')$ satisfies:
\begin{enumerate}[itemsep=2pt,topsep=4pt,leftmargin=16pt]
\item \textbf{Adjacency reflection: } $\Phi(\alpha,\beta)\iff \Phi'(f(\alpha),f(\beta))$.
\item \textbf{Label preservation: } $\lambda'\circ f=\lambda$.
\item \textbf{Equivariance: } $\forall u\in U\ \exists u'\in U'\ \text{s.t.}\ f\circ u=u'\circ f$.
\end{enumerate}
Items (1) and (2) imply $f$ is injective on vertices.

\subsubsection*{Atlas and Sub-Atlas}
\begin{itemize}[itemsep=2pt,topsep=2pt,leftmargin=16pt]
\item \textbf{Canonical full instance: } $V=R$, $E=\{\{\alpha,\beta\}:\langle\alpha,\beta\rangle=-1\}$, reconstructive $\lambda$, and $U=W(E_8)$.
\item \textbf{Induced sub-Atlas on } $I\subseteq\{1,\dots,240\}$: restrict $V,E,\lambda$; set $U=\{w|_V:w\in W(E_8),\ w(V)=V\}$.
\end{itemize}

\subsection{Example: The $E_7$ Subsystem}
Let
\[
V=\{\alpha\in R:\ \langle \alpha,\alpha_8\rangle=0\}\quad\text{(roots orthogonal to $\alpha_8$)},
\]
and $E=\{\{\alpha,\beta\}\subset V:\ \langle\alpha,\beta\rangle=-1\}$, with $\lambda$ restricted to $V$ and
\[
U=\{w|_V:\ w\in W(E_8),\ w(V)=V\}\cong W(E_7).
\]

\subsection{Automorphisms}
The $E_8$ Dynkin diagram has trivial diagram automorphism. The longest Weyl group element $w_0$ acts as $-\mathrm{Id}$ on $\mathbb{R}^8$, so $-\mathrm{Id}\in W(E_8)$. Hence
\[
\mathrm{Aut}(R,\Phi)=W(E_8).
\]

\subsection{Validation Checklist}
For implementation, verify:
\begin{enumerate}[itemsep=2pt,topsep=4pt,leftmargin=16pt]
\item \textbf{Cartan verification: } $G_{ij}=\langle\alpha_i,\alpha_j\rangle$ equals $A$.
\item \textbf{Root enumeration: } Enumerate $R$ from $E_8$; check $|R|=240$ and $\langle\alpha,\alpha\rangle=2$ for all $\alpha\in R$.
\item \textbf{Adjacency: } Build $\Phi$; every vertex has degree $56$.
\item \textbf{Weyl-invariance and transitivity: } $\Phi(\alpha,\beta)=\Phi(w\alpha,w\beta)$ for $w\in W(E_8)$ and $W(E_8)$ acts transitively on $R$.
\item \textbf{Label reconstruction: } For $V\subseteq R$, confirm $\mathrm{rec}\circ\lambda=\mathrm{id}_V$.
\end{enumerate}

% LaTeX section: Group Action, Orbit Decomposition, and Invariant Functionals
% Self-contained. Load standard packages in the parent document as needed.

% LaTeX section: Group Action, Orbit Decomposition, and Invariant Functionals
% Self-contained. Load standard packages in the parent document as needed.

\section{Group Action, Orbit Decomposition, and Invariant Functionals}

\subsection{Ambient Set and CRT Maps}
Let $P=\mathbb Z/48\mathbb Z$, $B=\mathbb Z/256\mathbb Z$, and $G=P\times B$ so $|G|=48\cdot 256=12{,}288$. Write the Chinese Remainder isomorphism for $P$ as
\[
\psi: P\to \mathbb Z/16\mathbb Z\times \mathbb Z/3\mathbb Z,\qquad \psi(p)=(a,c)=(p\bmod 16,\, p\bmod 3),
\]
with inverse (encoding)
\[
\varphi(a,c)\equiv 33a+16c\pmod{48},
\]
using $3^{-1}\equiv 11\pmod{16}$ and $16^{-1}\equiv 1\pmod{3}$.

Represent $a\in\mathbb Z/16\mathbb Z$ in binary as $a=a_0+2a_1+4a_2+8a_3$ with $a_i\in\{0,1\}$.

\subsection{Reference Subgroup and Action}
Let $V\cong (\mathbb Z/2\mathbb Z)^3$ act by flipping bits $a_0,a_1,a_2$ (bit $a_3$ fixed). Let $W\cong (\mathbb Z/2\mathbb Z)^8$ act on bytes by XOR. Define
\[
U_{\mathrm{ref}}=V\times W\cong (\mathbb Z/2\mathbb Z)^{11}.
\]
For $u=(v,w)\in U_{\mathrm{ref}}$ and $g=(p,b)\in G$:
\begin{enumerate}\itemsep2pt
  \item $(a,c)=\psi(p)$.
  \item $a' = (a\oplus v)$ on bits $a_0,a_1,a_2$ with $a_3$ unchanged.
  \item $b' = b\oplus w$.
  \item $p' = \varphi(a',c)$ and $u\cdot g=(p',b')$.
\end{enumerate}
The $11$ standard generators are commuting, each of order $2$, giving a simple undirected Cayley graph on $G$.

\subsection{Orbit Invariants, Freeness, and Tiling}
Define the orbit invariant
\[
\mathbf o(g)=(c, a_3)\in \mathbb Z/3\mathbb Z\times\{0,1\},\qquad c=p\bmod 3,\; a_3=\lfloor a/8\rfloor.
\]
\begin{proposition}
The action of $U_{\mathrm{ref}}$ on $G$ is free. The invariant $\mathbf o$ is preserved by $U_{\mathrm{ref}}$. Each orbit has size $2^{11}=2048$ and is determined by $(c,a_3)$. There are exactly six orbits, with representatives
\[
S=\{(\varphi(0,c),0),\, (\varphi(8,c),0): c\in\{0,1,2\}\}.
\]
Moreover, edges never cross orbits; each connected component of the Cayley graph equals a single orbit.
\end{proposition}
\begin{proof}[Sketch]
Nonzero $v$ changes $a$; nonzero $w$ changes $b$; hence stabilizers are trivial. The bits $a_0,a_1,a_2$ and the byte are toggled freely, while $c$ and $a_3$ are fixed, so orbits have size $2^{3+8}=2048$ and are indexed by $(c,a_3)\in\{0,1,2\}\times\{0,1\}$. The six anchors listed are in distinct orbits and tile $G$.
\end{proof}

\subsection{Cayley Graph and Flag Complex}
Let $\Gamma$ be the Cayley graph of $G$ with respect to the 11 generators. On each orbit, $\Gamma$ is the $11$-dimensional hypercube $Q^{11}$: per orbit $n=2^{11}=2048$ vertices and $m=11\cdot 2^{10}=11264$ edges. The entire graph has six connected components, totaling
\[
|V|=12{,}288,\qquad |E|=67{,}584.
\]
Hypercubes are triangle-free, thus the flag complex equals the $1$-skeleton. For a graph, $H_1$ is torsion-free and
\[
\beta_1 = m-n+\#\text{components}.
\]
Hence per orbit $\beta_1=11264-2048+1=9217$, and in total $\beta_1=6\cdot 9217=55{,}302$. All higher Betti numbers vanish.

\subsection{Schedule, Base Functional, and Invariance Scope}
Let $N=768$ and the canonical fair schedule
\[
\sigma(t)=(t\bmod 48,\; t\bmod 256),\qquad t=0,\dots,N-1.
\]
For an affine observable $f(p,b)=d+\alpha_p+\gamma_b$, define the base functional
\[
S_0(f)=\sum_{t=0}^{N-1} f(\sigma(t))=768\,d+16\sum_{p\in P}\!\alpha_p+3\sum_{b\in B}\!\gamma_b.
\]
\begin{proposition}
$S_0$ is invariant under arbitrary permutations of pages and bytes and under the $U_{\mathrm{ref}}$ action.
\end{proposition}
\begin{remark}[Limits for $k\ge 1$]
No polynomial weighting in $t$ alone equalizes arithmetic-progression sums induced by $\sigma$. In particular, for $\tau=t-(N-1)/2$, the partial sums $\sum_{t\equiv r\,(48)}\!\tau$ and $\sum_{t\equiv s\,(256)}\!\tau$ are residue-dependent, so $k\ge 1$ moment functionals are not invariant in general.
\end{remark}

\subsection{Permutation-Invariant Alternatives for First Moments}
Define integer-scaled per-residue weights
\[
\,\, w_p(t)=2\Big\lfloor\tfrac{t}{48}\Big\rfloor-15,\qquad w_b(t)=2\Big\lfloor\tfrac{t}{256}\Big\rfloor-2,
\]
so that for every page residue $r$ and byte residue $s$,
\[
\sum_{i=0}^{15} w_p(r+48i)=0,\qquad \sum_{i=0}^{2} w_b(s+256i)=0.
\]
Then the linear functionals
\[
S^{\mathrm{page}}_1(f)=\tfrac{1}{2}\sum_{t=0}^{N-1} w_p(t)\, f(\sigma(t)),\qquad
S^{\mathrm{byte}}_1(f)=\tfrac{1}{2}\sum_{t=0}^{N-1} w_b(t)\, f(\sigma(t))
\]
are invariant under arbitrary page permutations and byte permutations, respectively, and under the corresponding $U_{\mathrm{ref}}$ actions on $P$ and $B$.

\subsection{R96 Labeling}
Define $\lambda:B\to\{0,\dots,95\}$ by $\lambda(b)=b\bmod 96$. Multiplicities are $|\lambda^{-1}(r)|=3$ for $r\in\{0,\dots,63\}$ and $|\lambda^{-1}(r)|=2$ for $r\in\{64,\dots,95\}$, so $64\cdot 3+32\cdot 2=256$. Extend to $G$ by composition with the byte projection; no equivariance beyond Weyl-invariance is asserted.

\subsection{Operational Notes}
Record the orbit tag $\mathbf o(g)$ for indexing and caching. Since $\mathbf o$ is invariant, edges never cross orbits and each component can be processed independently. Sparse $\mathbb F_2$ or $\mathbb Z$ homology on the $67{,}584$-edge graph is straightforward; for graphs $\beta_1$ is identical over $\mathbb Z$ and $\mathbb F_2$.

\subsection{E8-Based 96-Vertex Resonance Graph}
Let $R\subset\mathbb R^8$ be the $E_8$ root system normalized by $\langle\alpha,\alpha\rangle=2$ for all $\alpha\in R$, so pairwise inner products lie in $\{-2,-1,0,1,2\}$. Fix an injective indexing map
\[
\iota:\{0,\dots,95\}\longrightarrow R,\qquad R_{96}:=\operatorname{im}(\iota),
\]
and define the induced simple graph $\Gamma_{E_8}^{(96)}$ by
\[
V(\Gamma_{E_8}^{(96)})=R_{96},\qquad \{\alpha,\beta\}\in E(\Gamma_{E_8}^{(96)})\iff (\alpha\neq\beta\text{ and }\langle\alpha,\beta\rangle=+1).
\]
No edges are added when $\langle\alpha,\beta\rangle\in\{0,1,2,-2\}$. This is the subgraph induced from the standard $E_8$ root graph under the $+1$ adjacency rule.

\paragraph{Adjacency conventions.} The full $E_8$ ResGraph elsewhere uses $-1$ adjacency (edges for $\langle\alpha,\beta\rangle=-1$). The 96-vertex graph here intentionally uses $+1$ adjacency. Both conventions are Weyl-equivariant. Note also that in this subsection all roots satisfy $\langle\alpha,\alpha\rangle=2$ by normalization.

\paragraph{Integration with the $R96$ labeling.} Use the existing byte partition to index $R_{96}$ via
\[
\lambda_{E_8}:B\to R_{96},\qquad \lambda_{E_8}(b):=\iota\big(b\bmod 96\big).
\]
Then $|\lambda_{E_8}^{-1}(\alpha)|=3$ for $\alpha\in\iota(\{0,\dots,63\})$ and $|\lambda_{E_8}^{-1}(\alpha)|=2$ for $\alpha\in\iota(\{64,\dots,95\})$, matching the multiplicities $64\cdot3+32\cdot2=256$. No compatibility with the $U_{\mathrm{ref}}$ action is required or claimed.

\paragraph{Independence from the Cayley layer.} $\Gamma_{E_8}^{(96)}$ lives solely on the label set and does not modify the Cayley graph $\Gamma$ or its flag complex. It is an auxiliary resonance structure replacing any ``96-dimensional'' phrasing with the precise term ``96-vertex subgraph of the $E_8$ root graph (adjacency $\langle\alpha,\beta\rangle=+1$)''.

\paragraph{Optional normalization.} One may choose $R_{96}$ closed under negation ($\alpha\in R_{96}\Rightarrow-\alpha\in R_{96}$); this adds no edges because $\langle\alpha,-\alpha\rangle=-2$, but can be convenient for symmetry bookkeeping. No further properties of $R_{96}$ are assumed.


\section{Atlas$\to E_8$ embedding: construction, invariants, uniqueness, and exact tests}

\subsection*{Construction}
Work in the type–II root subsystem of the $E_8$ lattice:
\[
x=\tfrac12(\pm1,\dots,\pm1)\in\Bbb R^8\quad\text{with an even number of minus signs}.
\]
Fix columns $h_1,\dots,h_7=e_i$, $h_8=\mathbf 1\in\Bbb F_2^7$ with $\sum_i h_i=0$.
For $u\in\Bbb F_2^7$ define signs $s_i(u)=(-1)^{u\cdot h_i}$ and
\[
x(u)=\tfrac12\big(s_1(u),\dots,s_8(u)\big).
\]
Select the $96$-subset
\[
S=\{u:u_7=0\}\ \cup\ \{u:u_7=1,\ u_1\oplus u_2\oplus u_3=0\},
\qquad
V=\{x(u):u\in S\}.
\]
Edges are declared by the $E_8$ inner product: $\langle x,y\rangle=1$.

\subsection*{Published data}
The 96-vector table and all invariants are provided as CSVs:
\texttt{verified\_E8\_96\_vectors.csv},
\texttt{verified\_allpair\_inner\_products.csv},
\texttt{verified\_edge\_inner\_products.csv},
\texttt{verified\_nonedge\_inner\_products.csv},
\texttt{verified\_degree\_distribution.csv},
\texttt{verified\_inner\_product\_counts.csv},
\texttt{verified\_nonedge\_inner\_product\_counts.csv}.

\subsection*{Verified invariants for the 96-set}
All computations use exact rationals.
\begin{itemize}\setlength\itemsep{2pt}
\item $|V|=96$, and $\langle x,x\rangle=2$ for all $x\in V$.
\item Distinct-pair inner-product multiset ($\binom{96}{2}=4560$):
\[
-2:~32,\qquad -1:~944,\qquad 0:~2544,\qquad +1:~1040.
\]
\item Graph statistics (edge iff $\langle x,y\rangle=+1$):
\[
\text{edges }=1040,\ \text{non-edges }=3520,\ 
\deg_{\min}=16,\ \deg_{\max}=25,\ \overline{\deg}=65/3.
\]
\end{itemize}
Antipodal pairs account for the $32$ occurrences of $-2$; this is consistent with $-I\in W(E_8)$.

\subsection*{Normalized $64$-set (no $-2$ among distinct pairs)}
Fix a canonical representative per antipodal class by requiring the first nonzero coordinate to be $+1/2$; then deduplicate.
This yields $V^{(64)}$ with $|V^{(64)}|=64$ and no antipodal collisions.
\begin{itemize}\setlength\itemsep{2pt}
\item Distinct-pair inner-product multiset ($\binom{64}{2}=2016$):
\[
-1:~224,\qquad 0:~1120,\qquad +1:~672.
\]
\item Graph is regular of degree $21$; edges $=672$, non-edges $=1344$.
\end{itemize}
Exports:
\texttt{normalized\_64\_vectors.csv},
\texttt{normalized\_64\_allpair\_inner\_products.csv},
\texttt{normalized\_64\_edge\_inner\_products.csv},
\texttt{normalized\_64\_nonedge\_inner\_products.csv},
\texttt{normalized\_64\_degree\_distribution.csv},
\texttt{normalized\_64\_inner\_product\_counts.csv},
\texttt{normalized\_64\_nonedge\_inner\_product\_counts.csv}.

\subsection*{Uniqueness up to $W(E_8)$}
The map $u\mapsto x(u)$ bijects $\Bbb F_2^7$ with even-parity type–II sign patterns.
The affine group $G=\mathrm{AGL}(7,2)=2^7\rtimes\mathrm{GL}(7,2)$ acts by $u\mapsto Au+t$,
inducing lattice automorphisms inside $W(E_8)$ and preserving inner products.
Our $96$-set has the form
\[
V=H_0\ \cup\ \{u\in H_1:m(u)=0\},\quad
H_b=\{u:\ell(u)=b\},
\]
with independent linear forms $\ell(u)=u_7$ and $m(u)=u_1\oplus u_2\oplus u_3$.
Since $\mathrm{GL}(7,2)$ acts transitively on ordered independent pairs $(\ell,m)$ and translations choose cosets, every
subset of the form “one full hyperplane plus half of a disjoint hyperplane’’ is $G$-equivalent; hence all such $96$-sets are $W(E_8)$-equivalent.
Presence or absence of antipodals is inconsequential because $-I\in W(E_8)$.

\subsection*{Reproducibility: exact-arithmetic test code}
\noindent\textbf{Inputs.} A CSV \texttt{verified\_E8\_96\_vectors.csv} with columns \texttt{x1,\dots,x8} containing strings in \{\texttt{"1/2"}, \texttt{"-1/2"}\}.
\medskip

\noindent\textbf{Python (exact checks and normalization).}
\begin{verbatim}
from fractions import Fraction
from itertools import combinations
import csv

# --- Load 96 vectors as Fractions ---
V = []
with open("verified_E8_96_vectors.csv", newline="") as f:
    rdr = csv.DictReader(f)
    cols = [f"x{i}" for i in range(1,9)]
    for row in rdr:
        v = [Fraction(1,2) if row[c].strip() == "1/2" else Fraction(-1,2) for c in cols]
        V.append(v)
assert len(V) == 96

# Convert to sign matrix S in {±1}
S = [[1 if a == Fraction(1,2) else -1 for a in v] for v in V]

# --- Exact checks for the 96-set ---
def ip(u, v):  # exact inner product in E8
    s = sum(uu*vv for uu, vv in zip(u, v))    # integer in {-8,-6,...,8}
    return Fraction(s, 4)

# Lengths
assert all(ip(s, s) == 2 for s in S)

# Pairwise histogram
from collections import Counter
hist96 = Counter(ip(S[i], S[j]) for i, j in combinations(range(96), 2))
# Expected: {-2:32, -1:944, 0:2544, +1:1040}
print("96-set IP histogram:", {int(k): v for k, v in hist96.items()})

# Degrees for edges defined by <x,y>=1
deg96 = [0]*96
for i, j in combinations(range(96), 2):
    if ip(S[i], S[j]) == 1:
        deg96[i] += 1; deg96[j] += 1
print("deg min/max/avg:", min(deg96), max(deg96), sum(deg96)/96)

# --- Normalize to 64-set: fix first coord positive, then deduplicate ---
def canon(row):
    return tuple(row if row[0] == 1 else (-x for x in row))

seen = set()
S64 = []
for r in S:
    c = canon(r)
    if c not in seen:
        seen.add(c)
        S64.append(list(c))
assert len(S64) == 64

# Verify 64-set invariants
hist64 = Counter(ip(S64[i], S64[j]) for i, j in combinations(range(64), 2))
assert Fraction(-2, 1) not in hist64  # no antipodal pairs
print("64-set IP histogram:", {int(k): v for k, v in hist64.items()})

deg64 = [0]*64
for i, j in combinations(range(64), 2):
    if ip(S64[i], S64[j]) == 1:
        deg64[i] += 1; deg64[j] += 1
print("64-set degree uniform:", min(deg64), max(deg64), sum(deg64)/64)
\end{verbatim}

\noindent\textbf{Expected output (96-set).}
\[
-2:32,\quad -1:944,\quad 0:2544,\quad +1:1040;\quad
\deg_{\min}=16,\ \deg_{\max}=25,\ \overline{\deg}=65/3.
\]

\noindent\textbf{Expected output (64-set).}
\[
-1:224,\quad 0:1120,\quad +1:672;\quad
\text{regular of degree }21.
\]


% =============================================================
\section{Restricted Initiality and Term Models for \textsc{ResGraph}}

\subsection{Signature, Axioms, and Models}
Let $\Sigma$ be the indexed total signature:
\begin{itemize}\itemsep2pt
  \item A binary product symbol $\otimes$ (disjoint orthogonal sum), with equations for commutativity, associativity, and unit; totality witnessed by indexing over orthogonal pairs.
  \item For each Dynkin automorphism $\Gamma$ of a simply-laced simple system, a unary symbol $\mathrm{fold}_{\Gamma}$.
  \item For each $I\subseteq\{1,\dots,8\}$, a unary symbol $\mathrm{Par}_I$ (parabolic restriction).
  \item For each admissible extension index $j$, a unary symbol $\mathrm{Aug}_j$ (inverse parabolic augmentation).
\end{itemize}
Add $96$ constants $(a_k)_{k\in A}$ pinning a marking $i:A\to X$. Include finite-limit axioms $\mathbf L$ enforcing:
\begin{enumerate}\itemsep2pt
  \item The equations characterizing $\otimes$, $\mathrm{fold}_{\Gamma}$, $\mathrm{Par}_I$, $\mathrm{Aug}_j$ and their compatibilities.
  \item The reflection-closure of $i(A)$ is a crystallographic root system of type $E_8$ with Cartan relations of $E_8$ (long-root norm $2$).
  \item Weyl-equivariance of these operations on the $E_8$-closure.
\end{enumerate}
A \emph{$\Sigma$-model with marking} (object of $\mathbf{Alg}_{96}$) is a model of the finite-limit sketch determined by $\Sigma$ and $\mathbf L$ plus the constants $(a_k)$.

\begin{proposition}[Basic properties]
$\mathbf{Alg}_{96}$ is locally presentable and monadic over \textbf{Set}. Hence there is a finitary monad $T$ on \textbf{Set} with $\mathbf{Alg}_{96}\simeq \mathbf{Set}^T$, and free (term) models exist.
\end{proposition}

\subsection{Free/Term Model}
Let $F(A)$ be the term model on the $96$-set $A$:
\[
F(A)=\mathsf{Tm}_\Sigma(A)/{\equiv_{\mathbf L}},\qquad \iota:A\to F(A),
\]
where $\equiv_{\mathbf L}$ is the least congruence generated by $\mathbf L$.

\begin{theorem}[Restricted Initiality, pointed form]\label{thm:restricted-initiality}
For every $\mathcal A\in \mathbf{Alg}_{96}$ interpreting $(a_k)$ as a marking $i:A\xrightarrow{\cong} S\subseteq \mathcal A$ whose reflection-closure is $E_8$, there exists a unique $\Sigma$-homomorphism
\[
\eta_{\mathcal A}:F(A)\to \mathcal A
\]
such that $\eta_{\mathcal A}\circ\iota=i$. The maps $\eta_{\mathcal A}$ are natural in $\mathcal A$. Hence $F(A)$ is initial in $\mathbf{Alg}_{96}$.
\end{theorem}
\begin{proof}[Proof sketch]
Define $\eta_{\mathcal A}$ by structural recursion on $\Sigma$-terms and quotient by $\equiv_{\mathbf L}$. Uniqueness by induction on term depth. Naturality from functoriality of term interpretation.
\end{proof}

\subsection{Unrestricted Initiality: Obstruction and Remedies}
Let $\mathrm{Conf}_{96}(\mathcal A)$ be the set of Atlas-type $96$-subsets of $\mathcal A$ whose reflection-closures are $E_8$. The Weyl group $W(E_8)$ acts; typically $\mathrm{Conf}_{96}(\mathcal A)$ is a (possibly empty) $W(E_8)$-torsor. A functorial map $F(A)\to\mathcal A$ without a chosen marking exists iff this torsor has a natural global section, which fails when automorphisms permute distinct configurations.

\begin{conjecture}[Unrestricted Initiality]
For every $\Sigma$-model $\mathcal A$ of $\mathbf L$ there is a morphism $F(A)\to\mathcal A$ unique up to automorphisms. Equivalently, every model carries a functorially natural Atlas-type $96$-configuration.
\end{conjecture}
\noindent\emph{Dependencies:} (C1) existence and Weyl-uniqueness of a $96$-configuration in every simply-laced rank-$8$ subsystem; (C2) preservation/detection by the indexed operations.

\paragraph{Forcing canonicity.}
(i) Add an orientation/height to select one configuration; (ii) move to groupoid semantics and claim initiality up to isomorphism; (iii) add a functorial selector $\sigma:\mathbf 1\Rightarrow \mathrm{Conf}_{96}$.

\subsection{Implementation Checklist}
\begin{enumerate}\itemsep2pt
  \item Enumerate index sets for $\mathrm{fold}_{\Gamma}$, $\mathrm{Par}_I$, $\mathrm{Aug}_j$ (totality).
  \item List the finite equations forcing the $E_8$ Cartan relations on the closure of $i(A)$.
  \item Prove $W(E_8)$-equivariance of $\mathrm{fold}_{\Gamma}$, $\mathrm{Aug}_j$, $\mathrm{Par}_I$ on the closure.
  \item Give a counterexample to unrestricted initiality: two disjoint $96$-configurations swapped by an automorphism.
\end{enumerate}
% =============================================================
% LaTeX section: Group Action, Orbit Decomposition, and Invariant Functionals
% Self-contained. Load standard packages in the parent document as needed.

% LaTeX section: Group Action, Orbit Decomposition, and Invariant Functionals
% Self-contained. Load standard packages in the parent document as needed.

\section{Group Action, Orbit Decomposition, and Invariant Functionals}

\subsection{Ambient Set and CRT Maps}
Let $P=\mathbb Z/48\mathbb Z$, $B=\mathbb Z/256\mathbb Z$, and $G=P\times B$ so $|G|=48\cdot 256=12{,}288$. Write the Chinese Remainder isomorphism for $P$ as
\[
\psi: P\to \mathbb Z/16\mathbb Z\times \mathbb Z/3\mathbb Z,\qquad \psi(p)=(a,c)=(p\bmod 16,\, p\bmod 3),
\]
with inverse (encoding)
\[
\varphi(a,c)\equiv 33a+16c\pmod{48},
\]
using $3^{-1}\equiv 11\pmod{16}$ and $16^{-1}\equiv 1\pmod{3}$.

Represent $a\in\mathbb Z/16\mathbb Z$ in binary as $a=a_0+2a_1+4a_2+8a_3$ with $a_i\in\{0,1\}$.

\subsection{Reference Subgroup and Action}
Let $V\cong (\mathbb Z/2\mathbb Z)^3$ act by flipping bits $a_0,a_1,a_2$ (bit $a_3$ fixed). Let $W\cong (\mathbb Z/2\mathbb Z)^8$ act on bytes by XOR. Define
\[
U_{\mathrm{ref}}=V\times W\cong (\mathbb Z/2\mathbb Z)^{11}.
\]
For $u=(v,w)\in U_{\mathrm{ref}}$ and $g=(p,b)\in G$:
\begin{enumerate}\itemsep2pt
  \item $(a,c)=\psi(p)$.
  \item $a' = (a\oplus v)$ on bits $a_0,a_1,a_2$ with $a_3$ unchanged.
  \item $b' = b\oplus w$.
  \item $p' = \varphi(a',c)$ and $u\cdot g=(p',b')$.
\end{enumerate}
The $11$ standard generators are commuting, each of order $2$, giving a simple undirected Cayley graph on $G$.

\subsection{Orbit Invariants, Freeness, and Tiling}
Define the orbit invariant
\[
\mathbf o(g)=(c, a_3)\in \mathbb Z/3\mathbb Z\times\{0,1\},\qquad c=p\bmod 3,\; a_3=\lfloor a/8\rfloor.
\]
\begin{proposition}
The action of $U_{\mathrm{ref}}$ on $G$ is free. The invariant $\mathbf o$ is preserved by $U_{\mathrm{ref}}$. Each orbit has size $2^{11}=2048$ and is determined by $(c,a_3)$. There are exactly six orbits, with representatives
\[
S=\{(\varphi(0,c),0),\, (\varphi(8,c),0): c\in\{0,1,2\}\}.
\]
Moreover, edges never cross orbits; each connected component of the Cayley graph equals a single orbit.
\end{proposition}
\begin{proof}[Sketch]
Nonzero $v$ changes $a$; nonzero $w$ changes $b$; hence stabilizers are trivial. The bits $a_0,a_1,a_2$ and the byte are toggled freely, while $c$ and $a_3$ are fixed, so orbits have size $2^{3+8}=2048$ and are indexed by $(c,a_3)\in\{0,1,2\}\times\{0,1\}$. The six anchors listed are in distinct orbits and tile $G$.
\end{proof}

\subsection{Cayley Graph and Flag Complex}
Let $\Gamma$ be the Cayley graph of $G$ with respect to the 11 generators. On each orbit, $\Gamma$ is the $11$-dimensional hypercube $Q^{11}$: per orbit $n=2^{11}=2048$ vertices and $m=11\cdot 2^{10}=11264$ edges. The entire graph has six connected components, totaling
\[
|V|=12{,}288,\qquad |E|=67{,}584.
\]
Hypercubes are triangle-free, thus the flag complex equals the $1$-skeleton. For a graph, $H_1$ is torsion-free and
\[
\beta_1 = m-n+\#\text{components}.
\]
Hence per orbit $\beta_1=11264-2048+1=9217$, and in total $\beta_1=6\cdot 9217=55{,}302$. All higher Betti numbers vanish.

\subsection{Schedule, Base Functional, and Invariance Scope}
Let $N=768$ and the canonical fair schedule
\[
\sigma(t)=(t\bmod 48,\; t\bmod 256),\qquad t=0,\dots,N-1.
\]
For an affine observable $f(p,b)=d+\alpha_p+\gamma_b$, define the base functional
\[
S_0(f)=\sum_{t=0}^{N-1} f(\sigma(t))=768\,d+16\sum_{p\in P}\!\alpha_p+3\sum_{b\in B}\!\gamma_b.
\]
\begin{proposition}
$S_0$ is invariant under arbitrary permutations of pages and bytes and under the $U_{\mathrm{ref}}$ action.
\end{proposition}
\begin{remark}[Limits for $k\ge 1$]
No polynomial weighting in $t$ alone equalizes arithmetic-progression sums induced by $\sigma$. In particular, for $\tau=t-(N-1)/2$, the partial sums $\sum_{t\equiv r\,(48)}\!\tau$ and $\sum_{t\equiv s\,(256)}\!\tau$ are residue-dependent, so $k\ge 1$ moment functionals are not invariant in general.
\end{remark}

\subsection{Permutation-Invariant Alternatives for First Moments}
Define integer-scaled per-residue weights
\[
\,\, w_p(t)=2\Big\lfloor\tfrac{t}{48}\Big\rfloor-15,\qquad w_b(t)=2\Big\lfloor\tfrac{t}{256}\Big\rfloor-2,
\]
so that for every page residue $r$ and byte residue $s$,
\[
\sum_{i=0}^{15} w_p(r+48i)=0,\qquad \sum_{i=0}^{2} w_b(s+256i)=0.
\]
Then the linear functionals
\[
S^{\mathrm{page}}_1(f)=\tfrac{1}{2}\sum_{t=0}^{N-1} w_p(t)\, f(\sigma(t)),\qquad
S^{\mathrm{byte}}_1(f)=\tfrac{1}{2}\sum_{t=0}^{N-1} w_b(t)\, f(\sigma(t))
\]
are invariant under arbitrary page permutations and byte permutations, respectively, and under the corresponding $U_{\mathrm{ref}}$ actions on $P$ and $B$.

\subsection{R96 Labeling}
Define $\lambda:B\to\{0,\dots,95\}$ by $\lambda(b)=b\bmod 96$. Multiplicities are $|\lambda^{-1}(r)|=3$ for $r\in\{0,\dots,63\}$ and $|\lambda^{-1}(r)|=2$ for $r\in\{64,\dots,95\}$, so $64\cdot 3+32\cdot 2=256$. Extend to $G$ by composition with the byte projection; no equivariance beyond Weyl-invariance is asserted.

\subsection{Operational Notes}
Record the orbit tag $\mathbf o(g)$ for indexing and caching. Since $\mathbf o$ is invariant, edges never cross orbits and each component can be processed independently. Sparse $\mathbb F_2$ or $\mathbb Z$ homology on the $67{,}584$-edge graph is straightforward; for graphs $\beta_1$ is identical over $\mathbb Z$ and $\mathbb F_2$.

\subsection{E8-Based 96-Vertex Resonance Graph}
Let $R\subset\mathbb R^8$ be the $E_8$ root system normalized by $\langle\alpha,\alpha\rangle=2$ for all $\alpha\in R$, so pairwise inner products lie in $\{-2,-1,0,1,2\}$. Fix an injective indexing map
\[
\iota:\{0,\dots,95\}\longrightarrow R,\qquad R_{96}:=\operatorname{im}(\iota),
\]
and define the induced simple graph $\Gamma_{E_8}^{(96)}$ by
\[
V(\Gamma_{E_8}^{(96)})=R_{96},\qquad \{\alpha,\beta\}\in E(\Gamma_{E_8}^{(96)})\iff (\alpha\neq\beta\text{ and }\langle\alpha,\beta\rangle=+1).
\]
No edges are added when $\langle\alpha,\beta\rangle\in\{0,-1,2,-2\}$. This is the subgraph induced from the standard $E_8$ root graph under the $+1$ adjacency rule.

\paragraph{Adjacency conventions.} The full $E_8$ ResGraph elsewhere uses $-1$ adjacency (edges for $\langle\alpha,\beta\rangle=-1$). The 96-vertex graph here intentionally uses $+1$ adjacency. Both conventions are Weyl-equivariant. Note also that in this subsection all roots satisfy $\langle\alpha,\alpha\rangle=2$ by normalization.

\paragraph{Integration with the $R96$ labeling.} Use the existing byte partition to index $R_{96}$ via
\[
\lambda_{E_8}:B\to R_{96},\qquad \lambda_{E_8}(b):=\iota\big(b\bmod 96\big).
\]
Then $|\lambda_{E_8}^{-1}(\alpha)|=3$ for $\alpha\in\iota(\{0,\dots,63\})$ and $|\lambda_{E_8}^{-1}(\alpha)|=2$ for $\alpha\in\iota(\{64,\dots,95\})$, matching the multiplicities $64\cdot3+32\cdot2=256$. No compatibility with the $U_{\mathrm{ref}}$ action is required or claimed.

\paragraph{Independence from the Cayley layer.} $\Gamma_{E_8}^{(96)}$ lives solely on the label set and does not modify the Cayley graph $\Gamma$ or its flag complex. It is an auxiliary resonance structure replacing any ``96-dimensional'' phrasing with the precise term ``96-vertex subgraph of the $E_8$ root graph (adjacency $\langle\alpha,\beta\rangle=+1$)''.

\paragraph{Optional normalization.} One may choose $R_{96}$ closed under negation ($\alpha\in R_{96}\Rightarrow-\alpha\in R_{96}$); this adds no edges because $\langle\alpha,-\alpha\rangle=-2$, but can be convenient for symmetry bookkeeping. No further properties of $R_{96}$ are assumed.

% =============================================================
\section{Proofs (mathematical core)}
\begin{description}
  \item[T1.] Product, quotient (fold), filtration, and augmentation in \textsc{ResGraph} output crystallographic root systems.
  \item[T2.] For each output: existence of a simple system; Cartan integrality; Dynkin identification; \emph{Weyl order = product of Coxeter exponents + 1}.
  \item[T3.] Restricted initiality: the free \textsc{ResGraph}-algebra on the Atlas--$E_8$ 96-seed is initial in the subcategory with a chosen 96-seed and laws $\mathbf L$.
\end{description}
Conjectural: unrestricted initiality (clearly labeled as conjecture elsewhere).

% =============================================================
\section{Computational evidence (tests, not proofs)}
\begin{description}
  \item[E1.] Atlas$\to E_8$ 96-set inner products over $\mathbb Q$: lengths $2$; pair multiset on \emph{unordered pairs with $\alpha\ne\beta$} is $\{-2{:}32,\,-1{:}944,\,0{:}2544,\,+1{:}1040\}$. Identity check: $32+944+2544+1040=4560=\binom{96}{2}$.
  \item[E2.] Edge graph (adjacency $\langle\alpha,\beta\rangle=+1$): $|E|=1040$; non-edges $\binom{96}{2}-1040=3520$; degree statistics $\min 16$, $\max 25$, average $\bar d=2|E|/|V|=65/3$.
  \item[E3.] \textsc{ResGraph} operations executed symbolically with rationals: extracted simple systems; Cartans integral with off-diagonal products in $\{0,1,2,3\}$; Dynkin types match the canonical list; Weyl orders match products of exponents $+1$.
\end{description}
Scope: E1--E3 are numerical certificates; they guide but do not replace proofs.

% =============================================================
\section{Artifact manifest}
\subsection*{Data (inputs, canonical CSVs)}
\begin{itemize}\itemsep2pt
  \item \texttt{verified\_E8\_96\_vectors.csv} --- 96 half-integer roots, fixed order.
  \item \texttt{iota\_96.csv} --- explicit indexing map $\iota$ (ordering of the 96 roots).
  \item \texttt{verified\_allpair\_inner\_products.csv} --- full $4560$ counts over $\{-2,-1,0,+1\}$.
  \item \texttt{verified\_edge\_inner\_products.csv} --- pairs with inner product $+1$.
  \item \texttt{verified\_nonedge\_inner\_products.csv} --- pairs with inner product in $\{-2,-1,0\}$.
  \item \texttt{verified\_degree\_distribution.csv} --- vertex degrees for the edge graph.
  \item \texttt{verified\_inner\_product\_counts.csv} --- global multiplicities.
  \item \texttt{verified\_nonedge\_inner\_product\_counts.csv} --- non-edge multiplicities.
  \item \texttt{adjacency\_plus1\_96.npz} --- sparse CSR of $\Gamma_{E_8}^{(96)}$ under $+1$ adjacency.
\end{itemize}
\subsection*{Derived (deterministic from the above)}
\begin{itemize}\itemsep2pt
  \item \texttt{cartan\_*.json} --- Cartan matrices per operation instance.
  \item \texttt{dynkin\_*.txt} --- recognized Dynkin type and rank.
  \item \texttt{weyl\_orders\_*.txt} --- Weyl-group orders.
  \item \texttt{resgraph\_ops\_spec.md} --- machine-readable op trees producing each output.
  \item \texttt{checks\_report.md} --- pass/fail per proof obligation.
\end{itemize}
\subsection*{Reproducibility}
Record SHA-256 for every file; fix rational arithmetic (no floats); fix ordering: lexicographic on vectors, then on pairs.

% =============================================================
\section{Formal-verification plan (Lean and Coq)}
\subsection*{Targets (proof obligations)}
\begin{description}
  \item[P0.] Exact model of $\mathbb Q^8$ with bilinear form and reflections.
  \item[P1.] 96-set correctness: each row has even parity; norm $2$.
  \item[P2.] Pairwise inner products in $\{-2,-1,0,+1\}$ with counts as in E1.
  \item[P3.] \textsc{ResGraph} operations preserve the ``crystallographic root system'' predicate.
  \item[P4.] For every output: existence of a simple system by constructive algorithm; Cartan integral with $a_{ii}=2$ and $a_{ij}a_{ji}\in\{0,1,2,3\}$.
  \item[P5.] Dynkin type equals the canonical type of the computed Cartan (decision procedure).
  \item[P6.] Weyl order equals product of Coxeter exponents $+1$ of that type.
  \item[P7.] Restricted initiality: existence and uniqueness (mod Weyl) of the \textsc{ResGraph}-algebra morphism from the free algebra on the 96-seed.
\end{description}
\subsection*{Lean (mathlib) roadmap}
Define
\verb|root_vec := vector ℚ 8|, an inner product \verb|ip|, reflections \verb|refl α x := x - 2*(ip x α / ip α α) • α|, and predicates
\verb|is_rootset|, \verb|is_crystallographic|, \verb|simple_system Π ⊆ X|, \verb|cartan Π|.
Avoid IO in proofs; re-derive CSV contents. Use \verb|linarith|, \verb|norm_num|, \verb|ring|, finite enumerations via \verb|finset|. Implement folds, filters, augments, Cartan$\to$Dynkin decision tree, and Weyl-order tables with proofs. Prove initiality by structural induction on syntax trees modulo $\mathbf L$.
\subsection*{Coq (QArith + MathComp) roadmap}
Mirror the Lean development using \verb|QArith| rationals and MathComp matrices; use small-scale reflection for finite decisions.
\subsection*{Tooling pins}
Record exact versions: Lean (version and mathlib SHA), Coq (version, QArith/MathComp SHAs).

% =============================================================
\section{Separation: proof vs. evidence}
Proof artifacts live in Lean/Coq files with completed theorems. Evidence artifacts live as CSV/MD plus reproducible generators. Reports label every claim as \textsc{PROVED} or \textsc{EVIDENCE} with pointers to files.

\appendix
\section{Machine check for G$_2$}\label{app:g2-code}
\begin{lstlisting}[language=Python,basicstyle=\ttfamily\small]
# G2 fold inside E8: minimal checks

from fractions import Fraction
import math

# Gram on span{ beta_l = a1+a2+a4, beta_s = a3 }:
# |beta_l|^2=6, |beta_s|^2=2, (beta_l,beta_s)=-3
G = [[Fraction(6,1), Fraction(-3,1)],
     [Fraction(-3,1), Fraction(2,1)]]

def dot(c,d):  # c,d in coefficient basis [beta_l, beta_s]
    return c[0]*(G[0][0]*d[0] + G[0][1]*d[1]) + c[1]*(G[1][0]*d[0] + G[1][1]*d[1])

def l2(c): return dot(c,c)

# simple roots in this rank-2 subsystem
aL = [Fraction(1,1), Fraction(0,1)]  # beta_l
aS = [Fraction(0,1), Fraction(1,1)]  # beta_s

# Cartan: A_ij = 2*(ai,aj)/(aj,aj)
def A(i,j):
    ai = aL if i==1 else aS
    aj = aL if j==1 else aS
    return 2*dot(ai,aj)/l2(aj)

A11, A12 = A(1,1), A(1,2)
A21, A22 = A(2,1), A(2,2)
print("Cartan:", [[int(A11), int(A12)],[int(A21), int(A22)]])

# angle and Coxeter m: theta = arccos( (aL,aS)/(|aL||aS|) )
cos_th = float(dot(aL,aS) / math.sqrt(float(l2(aL)*l2(aS))))
theta = math.acos(cos_th)                # in radians
m = round(math.pi / (math.pi - theta))   # since theta = pi - pi/m
print("cos(theta)=", f"{cos_th:.12f}")
print("theta_deg =", theta*180.0/math.pi)
print("m =", m)

# reflections and orbit
def refl(v, i):  # reflect v across ai
    ai = aL if i==1 else aS
    return [v[0] - 2*dot(v,ai)/l2(ai)*ai[0],
            v[1] - 2*dot(v,ai)/l2(ai)*ai[1]]

roots = set()
front = {tuple(aL), tuple(aS), tuple([-aL[0],-aL[1]]), tuple([-aS[0],-aS[1]])}
while front:
    new = set()
    for v in front:
        if v not in roots:
            roots.add(v)
            new.add(tuple(refl(list(v),1)))
            new.add(tuple(refl(list(v),2)))
    front = new

# count lengths
from collections import Counter
len_counts = Counter(l2(list(v)) for v in roots)
print("root_count =", len(roots))
print("length^2 classes:", {float(k):v for k,v in len_counts.items()})

# Weyl group size via matrices
def refl_matrix(i):
    I = [[Fraction(1,1),Fraction(0,1)],[Fraction(0,1),Fraction(1,1)]]
    row = G[i-1]
    if i==1: e = [Fraction(1,1), Fraction(0,1)]
    else:    e = [Fraction(0,1), Fraction(1,1)]
    M = [[(Fraction(2,1)/G[i-1][i-1]) * e[r] * row[c] for c in range(2)] for r in range(2)]
    return [[I[r][c] - M[r][c] for c in range(2)] for r in range(2)]

def matmul(A,B):
    return [[A[0][0]*B[0][0] + A[0][1]*B[1][0], A[0][0]*B[0][1] + A[0][1]*B[1][1]],
            [A[1][0]*B[0][0] + A[1][1]*B[1][0], A[1][0]*B[0][1] + A[1][1]*B[1][1]]]

def mateq(A,B):
    return all(A[r][c]==B[r][c] for r in range(2) for c in range(2))

S1, S2 = refl_matrix(1), refl_matrix(2)
group = [[[Fraction(1,1),Fraction(0,1)],[Fraction(0,1),Fraction(1,1)]]]
queue = [S1, S2]
while queue:
    M = queue.pop()
    if not any(mateq(M,N) for N in group):
        group.append(M)
        queue += [matmul(M,S1), matmul(M,S2), matmul(S1,M), matmul(S2,M)]
print("Weyl_order =", len(group))

# Expected:
# Cartan: [[2, -3], [-1, 2]]
# m = 6
# root_count = 12 with length^2 classes {2:6, 6:6}
# Weyl_order = 12
\end{lstlisting}

\section{Machine check for F$_4$}\label{sec:f4-fold}
\begin{lstlisting}[language=Python,basicstyle=\ttfamily\small]
# F4 via E6 quotient fold inside E8: minimal checks

from fractions import Fraction
import math
from collections import Counter

# Gram on {beta1,beta2,beta3,beta4}:
# |b1|^2=|b2|^2=1, |b3|^2=|b4|^2=2
# (b1,b2)=-1/2, (b2,b3)=-1, (b3,b4)=-1, others 0
G = [[Fraction(1,1),  Fraction(-1,2), Fraction(0,1),  Fraction(0,1)],
     [Fraction(-1,2), Fraction(1,1),  Fraction(-1,1), Fraction(0,1)],
     [Fraction(0,1),  Fraction(-1,1), Fraction(2,1),  Fraction(-1,1)],
     [Fraction(0,1),  Fraction(0,1),  Fraction(-1,1), Fraction(2,1)]]

n = 4
def dot(c,d): return sum(c[i]*G[i][j]*d[j] for i in range(n) for j in range(n))
def l2(c): return dot(c,c)

# simple roots in this folded basis
e = [[Fraction(1,1) if i==k else Fraction(0,1) for k in range(n)] for i in range(n)]

# Cartan A_ij = 2*(alpha_i,alpha_j)/(alpha_j,alpha_j)
A = [[2*dot(e[i],e[j])/l2(e[j]) for j in range(n)] for i in range(n)]
print("Cartan:")
for r in A: print([int(x) for x in r])

# Coxeter m_ij from product a_ij a_ji in {0,1,2,3}
def m_from_prod(p): return {0:2, 1:3, 2:4, 3:6}[p]
for (i,j) in [(0,1),(1,2),(2,3)]:
    p = A[i][j]*A[j][i]
    print(f"m_{{{i+1}{j+1}}} =", m_from_prod(int(p)))

# reflections and orbit
def refl(v,i): return [v[k] - 2*dot(v,e[i])/l2(e[i])*e[i][k] for k in range(n)]
roots, front = set(), set()
for i in range(n):
    front.add(tuple(e[i])); front.add(tuple([-x for x in e[i]]))
while front:
    new=set()
    for v in front:
        if v not in roots:
            roots.add(v)
            for i in range(n): new.add(tuple(refl(list(v),i)))
    front=new

cnt = Counter(l2(list(v)) for v in roots)
print("root_count =", len(roots))
print("length^2 classes:", {float(k):v for k,v in cnt.items()})

# Weyl group order from four reflections
def reflect_matrix(i):
    I = [[Fraction(1,1) if r==c else Fraction(0,1) for c in range(n)] for r in range(n)]
    S = [row[:] for row in I]
    row = G[i]
    for r in range(n):
        for c in range(n):
            S[r][c] -= (Fraction(2,1)/G[i][i])*(Fraction(1,1) if r==i else Fraction(0,1))*row[c]
    return S

def matmul(A,B):
    C=[[Fraction(0,1) for _ in range(n)] for _ in range(n)]
    for i in range(n):
        for k in range(n):
            for j in range(n):
                C[i][j]+=A[i][k]*B[k][j]
    return C

def mateq(A,B): return all(A[r][c]==B[r][c] for r in range(n) for c in range(n))

S=[reflect_matrix(i) for i in range(n)]
group=[[[Fraction(1,1) if r==c else Fraction(0,1) for c in range(n)] for r in range(n)]]
queue=S[:]
while queue:
    M=queue.pop()
    if not any(mateq(M,N) for N in group):
        group.append(M)
        for X in S:
            queue += [matmul(M,X), matmul(X,M)]
print("Weyl_order =", len(group))

# Expected:
# Cartan =
# [2,-1,0,0]; [-1,2,-1,0]; [0,-2,2,-1]; [0,0,-1,2]
# m_12=3, m_23=4, m_34=3
# root_count=48 with length^2 {1:24, 2:24}
# Weyl_order=1152
\end{lstlisting}


\section{Machine check E$_6$ (filtration)}\label{sec:e6-filtration}
\begin{lstlisting}[language=Python,basicstyle=\ttfamily\small]
# E6 via standard parabolic inside E8: minimal checks

from fractions import Fraction
from collections import Counter

# Bourbaki E6: trivalent at 4, edges:
# (1-3),(3-4),(4-5),(5-6), and arm (2-4)
n = 6
edges = {(1,3),(3,4),(4,5),(5,6),(2,4)}
edges |= {(j,i) for (i,j) in edges}

# Gram matrix: 2 on diag, -1 on edges, 0 else
G = [[Fraction(2,1) if i==j else
      (Fraction(-1,1) if (i+1,j+1) in edges else Fraction(0,1))
      for j in range(n)] for i in range(n)]

def dot(c,d): return sum(c[i]*G[i][j]*d[j] for i in range(n) for j in range(n))
def l2(c): return dot(c,c)

# simple roots
e = [[Fraction(1,1) if i==k else Fraction(0,1) for k in range(n)] for i in range(n)]

# Cartan matrix
A = [[2*dot(e[i],e[j])/l2(e[j]) for j in range(n)] for i in range(n)]
print("Cartan matrix A:")
for r in A: print([int(x) for x in r])

# det(A) and rank over Q
def det_rational(M):
    from copy import deepcopy
    M = [[Fraction(m,1) for m in row] for row in M]
    det = Fraction(1,1); m = len(M)
    for i in range(m):
        p = next((r for r in range(i,m) if M[r][i]!=0), None)
        if p is None: return Fraction(0,1)
        if p!=i: M[i],M[p]=M[p],M[i]; det*=-1
        det *= M[i][i]
        piv = M[i][i]
        for c in range(i,m): M[i][c] /= piv
        for r in range(i+1,m):
            f = M[r][i]
            if f!=0:
                for c in range(i,m): M[r][c] -= f*M[i][c]
    return det

def rank_rational(M):
    from copy import deepcopy
    M = [[Fraction(m,1) for m in row] for row in M]
    r=c=0; R=len(M); C=len(M[0])
    while r<R and c<C:
        p = next((i for i in range(r,R) if M[i][c]!=0), None)
        if p is None: c+=1; continue
        M[r],M[p]=M[p],M[r]; piv=M[r][c]
        for j in range(c,C): M[r][j]/=piv
        for i in range(R):
            if i!=r and M[i][c]!=0:
                f=M[i][c]
                for j in range(c,C): M[i][j]-=f*M[r][j]
        r+=1; c+=1
    return r

print("det(A) =", det_rational(A), "rank(A) =", rank_rational(A))

# Reflection closure
def refl(v,i):
    ai=e[i]
    return [v[k]-2*dot(v,ai)/l2(ai)*ai[k] for k in range(n)]

roots=set(); front=set()
for i in range(n):
    front.add(tuple(e[i])); front.add(tuple([-x for x in e[i]]))
while front:
    new=set()
    for v in front:
        if v not in roots:
            roots.add(v)
            for i in range(n): new.add(tuple(refl(list(v),i)))
    front=new

cnt = Counter(l2(list(v)) for v in roots)
print("root_count =", len(roots))
print("length^2 classes:", {float(k):v for k,v in cnt.items()})

# Weyl order from exponents
exponents = [1,4,5,7,8,11]
order = 1
for m in exponents: order *= (m+1)
print("|W(E6)| =", order)

# Expected:
# det(A)=3, rank=6
# root_count=72, all length^2=2
# |W(E6)|=51840
\end{lstlisting}

\section{Machine check E$_7$(augmentation)}\label{sec:e7-augmentation}
\begin{lstlisting}[language=Python,basicstyle=\ttfamily\small]
# E7 via standard parabolic inside E8 with augmentation at the junction

from fractions import Fraction
from collections import Counter

# Bourbaki E7 inside E8: trivalent at node 4
# edges: (1-3),(3-4),(4-5),(5-6),(6-7) and arm (2-4)
n = 7
edges = {(1,3),(3,4),(4,5),(5,6),(6,7),(2,4)}
edges |= {(j,i) for (i,j) in edges}

# Gram matrix: 2 on diag, -1 on edges, 0 else
G = [[Fraction(2,1) if i==j else
      (Fraction(-1,1) if (i+1,j+1) in edges else Fraction(0,1))
      for j in range(n)] for i in range(n)]

def dot(c,d): return sum(c[i]*G[i][j]*d[j] for i in range(n) for j in range(n))
def l2(c): return dot(c,c)

# simple roots basis
e = [[Fraction(1,1) if i==k else Fraction(0,1) for k in range(n)] for i in range(n)]

# Cartan matrix
A = [[2*dot(e[i],e[j])/l2(e[j]) for j in range(n)] for i in range(n)]
print("Cartan matrix A:")
for r in A: print([int(x) for x in r])

# det(A) and rank over Q
def det_rational(M):
    M = [[Fraction(m,1) for m in row] for row in M]
    det = Fraction(1,1); m = len(M)
    for i in range(m):
        p = next((r for r in range(i,m) if M[r][i]!=0), None)
        if p is None: return Fraction(0,1)
        if p!=i: M[i],M[p]=M[p],M[i]; det*=-1
        piv = M[i][i]; det *= piv
        for c in range(i,m): M[i][c] /= piv
        for r in range(i+1,m):
            f = M[r][i]
            if f!=0:
                for c in range(i,m): M[r][c] -= f*M[i][c]
    return det

def rank_rational(M):
    M = [[Fraction(m,1) for m in row] for row in M]
    r=c=0; R=len(M); C=len(M[0])
    while r<R and c<C:
        p = next((i for i in range(r,R) if M[i][c]!=0), None)
        if p is None: c+=1; continue
        M[r],M[p]=M[p],M[r]; piv=M[r][c]
        for j in range(c,C): M[r][j]/=piv
        for i in range(R):
            if i!=r and M[i][c]!=0:
                f=M[i][c]
                for j in range(c,C): M[i][j]-=f*M[r][j]
        r+=1; c+=1
    return r

print("det(A) =", det_rational(A), "rank(A) =", rank_rational(A))

# Reflection-closure
def refl(v,i):
    ai=e[i]
    return [v[k]-2*dot(v,ai)/l2(ai)*ai[k] for k in range(n)]

roots,setfront=set(),set()
front=set()
for i in range(n):
    front.add(tuple(e[i]))
    front.add(tuple([-x for x in e[i]]))

while front:
    new=set()
    for v in front:
        if v not in roots:
            roots.add(v)
            for i in range(n):
                new.add(tuple(refl(list(v),i)))
    front=new

cnt = Counter(l2(list(v)) for v in roots)
print("root_count =", len(roots))
print("length^2 classes:", {float(k): v for k,v in cnt.items()})

# Weyl order from exponents
exponents = [1,5,7,9,11,13,17]
order = 1
for m in exponents: order *= (m+1)
print("|W(E7)| =", order)

# Expected:
# det(A)=2, rank=7
# root_count=126, all length^2=2
# |W(E7)|=2,903,040
\end{lstlisting}

\nocite{*}
\bibliographystyle{unsrt}
\bibliography{references}
\end{document}


